\chapter{Речовини, які не задають тип нейрона}
\label{sec:othermediators}

У розділі~\ref{sec:neurontypes} ми розсортували нейрони за головним медіатором і отримали десять типів (таблиця~\ref{tab:types}). Тепер час додати ускладнення. Нейрон викидає не тільки головний медіатор, а в мозку, крім головних медіаторів, працюють й інші речовини. Вони змінюють те, як мережа передає і зберігає сигнали.

Крім адресної та широкомовної шин (підрозділ~\ref{sec:buses}) є й третя: \emph{зворотний канал}. Кілька речовин викидає не відправник, а \emph{отримувач}, і вони йдуть через щілину назад, до виходу відправника, і змінюють, скільки медіатора той викине наступного разу. У мережах зв'язку це схоже на підтвердження приймання, яке відправник використовує, щоб налаштувати свою передачу. Саме цим каналом нейрони користуються, коли змінюють ваги своїх зв'язків: ним вихід відправника дізнається про активність отримувача, тобто про другий множник правил навчання розділу~\ref{sec:dofa}.

Повний список відомих медіаторів довгий: понад сотня речовин, більшість із них --- короткі білкові ланцюжки, нейропептиди. Але всі вони вкладаються в кілька класів, і всередині кожного класу речовини поводяться схоже. Речовини, які не бувають головним медіатором, зібрано в таблиці~\ref{tab:othermediators}, з тими самими трьома питаннями, які поставив би інженер: якою шиною йде сигнал, чи швидко він діє і який у нього зазвичай знак. Тип нейрона вони не задають: їх викидають разом із головним медіатором, або їх викидають не нейрони, або вони йдуть у зворотний бік.

\begin{table}[t]
\centering
\footnotesize
\setlength{\tabcolsep}{4pt}
\renewcommand{\arraystretch}{1.12}
\begin{tabular}{>{\raggedright}p{2.25cm}>{\raggedright}p{2.8cm}>{\raggedright}p{1.8cm}>{\raggedright}p{1.5cm}>{\centering}p{0.8cm}>{\raggedright\arraybackslash}p{4.3cm}}
\toprule
Клас & Речовина & Шина & Швид\-кість & Знак & Роль в обчисленні \\
\midrule
Амінокислоти & D-серин & локальна & повільно & І & другий ключ для рецептора глутамату: умова «І» \\
\midrule
Моноаміни & слідові аміни & широко\-мовна & повільно & $\pm$ & тонке налаштування моноамінових каналів \\
\midrule
\multirow{2}{2.25cm}{Пурини}
 & АТФ & адресна & швидко і повільно & $+$ & супутній медіатор; сигнал між нейронами і глією \\
 & аденозин & об'ємна & повільно & $-$ & накопичується під час активності: лічильник навантаження~\cite{Dunwiddie2001} \\
\midrule
\multirow{8}{2.25cm}{Нейропептиди (понад сотню)}
 & енкефаліни, ендорфіни, динорфін & об'ємна & повільно & $-$ & придушення передачі \\
 & речовина P & об'ємна & повільно & $+$ & повільне підсилення \\
 & нейропептид Y & об'ємна & повільно & $-$ & повільне гальмування \\
 & соматостатин & об'ємна & повільно & $-$ & повільне гальмування, супроводжує GABA \\
 & вазоактивний інтестинальний пептид (ВІП) & об'ємна & повільно & $+$ & розгальмовування: гальмування гальмівних нейронів \\
 & холецистокінін & об'ємна & повільно & $\pm$ & супроводжує GABA в частині гальмівних нейронів \\
 & орексин & широко\-мовна & повільно & $+$ & стабілізатор режиму неспання \\
 & окситоцин, вазопресин & широко\-мовна & повільно & $\pm$ & глобальні налаштування поведінки \\
\midrule
\multirow{2}{2.25cm}{Гази}
 & оксид азоту NO & зворотна, об'ємна & повільно & $\pm$ & проходить крізь мембрани; зворотний сигнал під час зміни ваг~\cite{Garthwaite2008} \\
 & CO, H$_2$S & об'ємна & повільно & $\pm$ & роль вивчено мало \\
\midrule
Ліпіди & ендоканабіноїди (анандамід, 2-АГ) & зворотна & повільно & $-$ & отримувач зменшує викид у відправника: зворотний зв'язок за вагою~\cite{WilsonNicoll2001} \\
\bottomrule
\end{tabular}
\caption{Речовини, які не задають тип нейрона. \emph{Шина}, \emph{швидкість} і \emph{знак} --- як у таблиці~\ref{tab:types}; крім того, об'ємна шина --- медіатор розпливається в міжклітинному просторі навколо місця викиду, зворотна --- йде від отримувача назад до відправника, локальна --- діє поруч із місцем викиду. «І» --- дозвільний сигнал: сам по собі струму не викликає, але без нього не спрацьовує інший вхід, як другий вхід логічного елемента «І». Нейропептиди майже завжди викидаються разом із головним медіатором і здебільшого за високої частоти імпульсів~\cite{vandenPol2012}.}
\label{tab:othermediators}
\end{table}
